\documentclass[11pt]{article}
\usepackage{amsmath,amssymb,amsthm}
\usepackage[margin=1in]{geometry}
\newtheorem{theorem}{Theorem}
\newtheorem{corollary}[theorem]{Corollary}
\title{The intersection property in rank three for simple groups}
\author{Working mathematical note}
\date{}
\begin{document}
\maketitle

This note supplies an elementary proof of the intersection assertion used by
the subgroup counting route for the unresolved rank-three cases of L026.
The assertion is treated as a known consequence of simplicity; no claim of
historical novelty is made. The definitions and the correspondence with
regular polytopes are those recalled in Section~2 of~\cite{LM2022}.
This note establishes no new numerical count or open conjecture, and it is
not a Lean formalization.

\begin{theorem}\label{thm:simple}
Let $G$ be a finite nonabelian simple group, and suppose that
\[
 G=\langle a,b,c\rangle,\qquad
 a^2=b^2=c^2=1,\qquad a,b,c\ne1,\qquad a\ne c,\qquad ac=ca.
\]
Then $(G;a,b,c)$ has the full intersection property: for all
$A,B\subseteq\{0,1,2\}$, with $r_0=a,r_1=b,r_2=c$,
\[
 \langle r_i:i\in A\rangle\cap\langle r_i:i\in B\rangle
       =\langle r_i:i\in A\cap B\rangle.
\]
In particular, it is a string C-group representation of rank three.
\end{theorem}

\begin{proof}
Put
\[
 D_1=\langle a,b\rangle,\quad D_2=\langle b,c\rangle,\quad
 R_1=\langle ab\rangle,\quad R_2=\langle bc\rangle.
\]
Neither adjacent pair commutes. Indeed, if $a$ commuted with $b$, it would
commute with all three generators and belong to $Z(G)$, which is trivial.
The argument for $b,c$ is the same, using $c$. Thus $ab$ and $bc$ have order
greater than two. Each $D_i$ is a dihedral group, $R_i$ is its cyclic rotation
subgroup of index two, and $b$ is a reflection in each $D_i$.

Let $I=D_1\cap D_2$ and $J_i=I\cap R_i$. Since $b\in I\setminus R_i$,
the restriction to $I$ of $D_i\longrightarrow D_i/R_i$ is onto. Consequently
$J_i$ is cyclic of index two in $I$, and
\[
                         I=\langle J_i,b\rangle.
\]
Conjugation by $b$ inverts $J_i$, so this is again a dihedral presentation.

Suppose first that $|I|>4$. Then $|J_i|>2$. A dihedral group whose rotation
subgroup has order greater than two has a unique cyclic subgroup of index
two. To see this directly, every element outside the rotation subgroup has
order two; a generator of a cyclic subgroup of that larger order must
therefore lie in the rotation subgroup, and equality follows from orders.
It follows that $J_1=J_2=:J$.

The normality required here does not require $I$ to be normal in either
$D_i$. Each $J_i$ is a subgroup of the cyclic group $R_i$, hence is
characteristic in $R_i$. Since $R_i\lhd D_i$, we have $J_i\lhd D_i$.
Thus $J$ is normalized by both $D_1$ and $D_2$, which together generate $G$.
It is a nontrivial cyclic normal subgroup of $G$, contrary to simplicity
and the assumption that $G$ is nonabelian.

Suppose next that $|I|=4$. Then $I$ is a Klein four group. Write
$J_i=\langle z_i\rangle$. An involution in the rotation subgroup of a
dihedral group is central, so $z_i\in Z(D_i)$. If $z_1=z_2$, this common
nonidentity element centralizes both $D_1$ and $D_2$, and hence all of $G$,
again a contradiction.

Otherwise $z_1,z_2$ are distinct nonidentity elements of $I$, and neither
equals $b$. Therefore $b=z_1z_2$. Set
\[
                 H=\langle a,z_2\rangle,\qquad
                 K=\langle c,z_1\rangle.
\]
All four cross-generator pairs commute: $a$ commutes with $c$ by hypothesis
and with $z_1$ because $z_1\in Z(D_1)$; similarly $z_2$ commutes with $c$;
and $z_1,z_2$ lie in the abelian group $I$. Thus $[H,K]=1$.
The product $HK$ contains $a,c$ and $b=z_1z_2$, so $G=HK$.
It follows that $H\lhd G$. Since $a\ne1$, simplicity implies $H=G$.
But then $K\leq Z(G)=1$, contradicting $c\in K\setminus\{1\}$.
This argument permits degenerate auxiliary factors $H,K$; their precise
dihedral orders are irrelevant.

The only remaining possibility is
\begin{equation}\label{eq:middle}
                         D_1\cap D_2=\langle b\rangle.
\end{equation}
Now $V=\langle a,c\rangle$ is a Klein four group. Its intersection with
$D_1$ contains $\langle a\rangle$. A larger subgroup of $V$ containing
$a$ would be $V$ itself, putting $c$ in $D_1$ and forcing $G=D_1$.
A finite dihedral group is not nonabelian simple. Hence
\begin{equation}\label{eq:ends}
                 D_1\cap V=\langle a\rangle,
                 \qquad D_2\cap V=\langle c\rangle.
\end{equation}

For completeness, these identities give the whole intersection property,
not just the displayed maximal-parabolic cases. The three generators are
distinct: an adjacent equality would put $G$ in the group generated by the
other two involutions. No generator belongs to the parabolic generated by
the other two, since that would make $G$ dihedral or equal to $V$.
Consequently a singleton parabolic meets the opposite two-generator
parabolic trivially, and distinct singleton parabolics meet trivially.
The intersections of distinct two-generator parabolics are exactly
\eqref{eq:middle} and~\eqref{eq:ends}. Nested subsets, the empty subset and
the full subset give their asserted intersections immediately. These cases
exhaust all pairs of subsets of $\{0,1,2\}$.
\end{proof}

\begin{corollary}
For a finite nonabelian simple group $G$, let $e(G)$ be the number of labelled
generating triples of nonidentity involutions with distinct commuting first
and third entries. Then $e(G)$ counts precisely its labelled rank-three
string C-group representations. The number up to ordered isomorphism is
\[
                          \frac{e(G)}{|\operatorname{Aut}G|}.
\]
\end{corollary}
\begin{proof}
The theorem proves the previously additional intersection requirement.
Conversely, each rank-three string C-group representation supplies such a
triple. An automorphism fixing all three entries fixes their generated group
pointwise and is therefore the identity. The automorphism action is free,
and each orbit has size $|\operatorname{Aut}G|$.
\end{proof}

This corollary concerns ordered isomorphism classes. Counting also up to
duality still requires the fixed classes of the reversal $(a,b,c)\mapsto
(c,b,a)$. That separate count is not inferred here. Nor does the theorem
permit imposing the intersection property on the proper-subgroup terms of
Hall inversion: those terms must count all generating triples in the
original domain.

\begin{thebibliography}{1}
\bibitem{LM2022}
D.~Leemans and J.~Mulpas,
\emph{The string C-group representations of sporadic groups},
Art Discrete Appl. Math. \textbf{5} (2022), \#P3.09, Section~2.
\end{thebibliography}
\end{document}
